% =============================================================
% บทที่ 1 - บทนำ
% เนื้อหาแปลง/เรียบเรียงจาก Progress/Proposal Project1.docx ฉบับปรับปรุงล่าสุด
% (หัวข้อ 2 หลักการและเหตุผล, 3 วัตถุประสงค์, 6 ขอบเขตและข้อจำกัดของการวิจัย,
% 8 ประโยชน์ที่คาดว่าจะได้รับ)
% =============================================================

% -----------------------------------------------------------
\section{ความเป็นมาและความสำคัญของปัญหา}
% -----------------------------------------------------------

จังหวัดขอนแก่นนับเป็นจังหวัดที่มีศักยภาพทางการท่องเที่ยวสูง ด้วยความหลากหลายของแหล่งท่องเที่ยว ทั้งทางธรรมชาติ วัฒนธรรม ประวัติศาสตร์ ตลอดจนแหล่งท่องเที่ยวเชิงวิชาการที่โดดเด่นอย่าง ``ไดโนเสาร์'' ซึ่งถือเป็นอัตลักษณ์ที่สร้างชื่อเสียงให้กับจังหวัด ขอนแก่นจึงเป็นหนึ่งในจุดหมายสำคัญของนักท่องเที่ยวทั้งชาวไทยและต่างชาติที่ต้องการสัมผัสประสบการณ์ที่หลากหลายและแตกต่าง

อย่างไรก็ตาม ปัญหาสำคัญที่ยังคงปรากฏอยู่คือ การเข้าถึงข้อมูลการท่องเที่ยวที่กระจัดกระจายอยู่ในหลายช่องทาง ไม่ว่าจะเป็นเว็บไซต์ หน่วยงานราชการ เพจโซเชียลมีเดีย หรือแอปพลิเคชันทั่วไป นักท่องเที่ยวจำเป็นต้องใช้เวลาค้นหาข้อมูลจากหลายแหล่ง ซึ่งมักพบปัญหาว่าข้อมูลไม่ครบถ้วน ไม่ทันสมัย หรือขาดการรองรับภาษาสำหรับนักท่องเที่ยวต่างชาติ ส่งผลให้การวางแผนท่องเที่ยวไม่มีประสิทธิภาพและอาจทำให้นักท่องเที่ยวพลาดโอกาสในการเข้าถึงประสบการณ์ที่ควรได้รับ

เพื่อแก้ไขข้อจำกัดดังกล่าว จึงเกิดแนวคิดในการพัฒนาโครงงาน ``ไดโน: เว็บแอปพลิเคชันศูนย์กลางของการท่องเที่ยวจังหวัดขอนแก่น'' ซึ่งเป็นแอปพลิเคชันที่มุ่งเน้นการรวมศูนย์ข้อมูลการท่องเที่ยวอย่างครบวงจร ภายในแอปประกอบด้วยฟังก์ชันการค้นหาแหล่งท่องเที่ยว ร้านอาหาร ที่พัก และกิจกรรมต่าง ๆ ระบบวางแผนการเดินทางที่จัดลำดับสถานที่ให้สอดคล้องกับระยะเวลา ความสนใจ และงบประมาณของผู้ใช้ รวมถึงระบบ Chatbot ที่ตอบคำถามโดยอ้างอิงเฉพาะข้อมูลในฐานความรู้ของระบบ ช่วยให้การค้นหาข้อมูลและการวางแผนทริปเป็นไปอย่างสะดวก รวดเร็ว และทันสมัย

นอกจากการอำนวยความสะดวกแก่นักท่องเที่ยวแล้ว แอปพลิเคชันยังเป็นเครื่องมือสำคัญในการสนับสนุนผู้ประกอบการท้องถิ่นให้สามารถโปรโมทธุรกิจได้ตรงกลุ่มเป้าหมาย อีกทั้งยังสามารถเก็บสถิติการใช้งาน เช่น จุดท่องเที่ยวยอดนิยม ช่วงฤดูกาลท่องเที่ยว หรือกิจกรรมที่ได้รับความสนใจสูง ซึ่งข้อมูลเหล่านี้สามารถนำไปใช้วิเคราะห์แนวโน้มและวางแผนการพัฒนาการท่องเที่ยวของจังหวัดในอนาคต

ดังนั้น การพัฒนาแอปพลิเคชัน ``ไดโน'' จึงไม่เพียงเป็นการแก้ปัญหาด้านการเข้าถึงข้อมูลการท่องเที่ยวเท่านั้น แต่ยังมีบทบาทสำคัญในการสร้างประสบการณ์การท่องเที่ยวที่ครบถ้วนสำหรับผู้มาเยือน และสนับสนุนเศรษฐกิจท้องถิ่นอย่างยั่งยืน ทั้งยังเป็นการวางรากฐานที่มั่นคงต่อการพัฒนาการท่องเที่ยวของจังหวัดขอนแก่นในระยะยาว

% -----------------------------------------------------------
\section{วัตถุประสงค์}
% -----------------------------------------------------------

\begin{enumerate}[label=\arabic*)]
    \item เพื่อพัฒนาเว็บแอปพลิเคชันรวมศูนย์ข้อมูลการท่องเที่ยวของจังหวัดขอนแก่น ที่นักท่องเที่ยวสามารถเข้าถึงข้อมูลแหล่งท่องเที่ยว ร้านอาหาร และกิจกรรมต่าง ๆ ที่เกิดขึ้นในจังหวัดขอนแก่น
    \item เพื่อพัฒนาและประยุกต์ใช้ระบบตอบโต้ผ่านโปรแกรมอัตโนมัติ (Chatbot) ในการให้บริการข้อมูลการท่องเที่ยว ตอบข้อซักถาม และให้คำแนะนำแก่ผู้ใช้งาน
    \item เพื่อพัฒนาระบบวางแผนการท่องเที่ยวที่สามารถจัดลำดับแผนการเดินทางให้สอดคล้องกับความต้องการของผู้ใช้
\end{enumerate}

% -----------------------------------------------------------
\section{ขอบเขตและข้อจำกัดของระบบ}
% -----------------------------------------------------------

\subsection{ขอบเขต}

การวิจัยนี้มุ่งเน้นการพัฒนาระบบวางแผนการเดินทาง (Trip Planning System) สำหรับจังหวัดขอนแก่น โดยประยุกต์ใช้เทคโนโลยีโมเดลภาษาขนาดใหญ่ (Large Language Models: LLMs) ร่วมกับสถาปัตยกรรม Retrieval-Augmented Generation (RAG) เพื่อสร้างแผนการเดินทางที่สอดคล้องกับความต้องการของผู้ใช้ และสามารถอ้างอิงข้อมูลจากฐานความรู้ที่ได้รับการจัดเก็บและปรับปรุงอย่างต่อเนื่อง ทั้งนี้ ระบบต้นแบบบางส่วนของงานวิจัยนี้ได้รับการพัฒนาและทดสอบการทำงานเบื้องต้นแล้ว ณ ช่วงเวลาที่จัดทำเค้าโครงโครงงานนี้

ระบบที่พัฒนาประกอบด้วยองค์ประกอบหลัก ดังนี้

\begin{enumerate}[label=(\arabic*)]
\item ระบบวางแผนการเดินทาง (Trip Planning System) ใช้เทคโนโลยี LLMs ร่วมกับ RAG ในการสร้างแผนการเดินทาง โดยพิจารณาจากข้อมูลที่ผู้ใช้กำหนด ได้แก่ วันที่และระยะเวลาเดินทาง จังหวะการท่องเที่ยว (Pace) ซึ่งส่งผลต่อจำนวนสถานที่ที่บรรจุในแผนแต่ละวัน ความสนใจ (เช่น วัด พิพิธภัณฑ์ ธรรมชาติ หรือร้านอาหาร) งบประมาณในการเดินทาง และสถานที่ที่ผู้ใช้ต้องการไปเป็นพิเศษ (Must-go) เพื่อสร้างแผนการเดินทางที่เหมาะสมและสอดคล้องกับข้อกำหนดของผู้ใช้ ทั้งนี้ ระบบมีกลไกตรวจสอบข้อจำกัดพื้นฐาน (Constraint Verification) ที่ทำงานอิสระจากแบบจำลองภาษา เพื่อยืนยันความสมเหตุสมผลของแผนก่อนนำเสนอผลลัพธ์แก่ผู้ใช้
\item ระบบฐานความรู้และการสืบค้นข้อมูล (Knowledge Retrieval System) ใช้ฐานข้อมูล PostgreSQL ที่มีส่วนขยายสำหรับจัดเก็บและค้นคืนข้อมูลเชิงเวกเตอร์ (pgvector) เป็นแหล่งจัดเก็บข้อมูลแหล่งท่องเที่ยว ร้านอาหาร ที่พัก กิจกรรม และเหตุการณ์ (Events) ภายในจังหวัดขอนแก่นแบบรวมศูนย์ โดยข้อมูลหลักรวบรวมจาก Google Places API และข้อมูลที่ผู้ดูแลระบบเพิ่มเข้าสู่ระบบผ่านหน้าจอผู้ดูแล ทั้งนี้ ในระยะถัดไปของงานวิจัยมีแผนขยายแหล่งข้อมูลให้ครอบคลุมแหล่งข้อมูลสาธารณะอื่นเพิ่มเติม เช่น สื่อสังคมออนไลน์ของผู้ประกอบการในพื้นที่ เพื่อเพิ่มความครบถ้วนและความทันสมัยของฐานความรู้ ข้อมูลที่จัดเก็บนี้สนับสนุนการค้นคืนข้อมูลสำหรับระบบ RAG และระบบตอบคำถามอัตโนมัติ
\item ระบบ Chatbot (Chatbot System) ใช้เทคโนโลยี LLMs และ RAG ในการตอบคำถามเกี่ยวกับการท่องเที่ยวจังหวัดขอนแก่น โดยอ้างอิงเฉพาะข้อมูลที่ค้นคืนได้จากฐานความรู้ของระบบ พร้อมกลไกตรวจจับกรณีที่ไม่พบข้อมูลตรงกับคำถาม เพื่อลดโอกาสการตอบคำถามที่คลาดเคลื่อนจากข้อมูลจริง (Hallucination) เพื่อให้ข้อมูลที่เกี่ยวข้องกับสถานที่ท่องเที่ยว ร้านอาหาร กิจกรรม การเดินทาง และข้อมูลที่ผู้ดูแลระบบเผยแพร่
\end{enumerate}

เว็บแอปพลิเคชัน (Web-based Application) รองรับการใช้งานของผู้ใช้ 2 กลุ่ม ได้แก่

\begin{enumerate}[label=(\arabic*)]
\item ผู้ใช้งานทั่วไป สามารถค้นหาข้อมูลแหล่งท่องเที่ยว กิจกรรม และสร้างแผนการเดินทางผ่านเว็บไซต์ รวมถึงสมัครสมาชิกและเข้าสู่ระบบเพื่อใช้งานระบบสะสมคะแนนจากการสแกน QR Code ณ สถานที่ที่เข้าร่วมโครงการ ซึ่งเป็นความสามารถเพิ่มเติม (Optional Feature) ของระบบ
\item ผู้ดูแลระบบ (Administrator) สามารถจัดการข้อมูลแหล่งท่องเที่ยว ร้านอาหาร กิจกรรมหรืออีเวนต์ในจังหวัดขอนแก่น รวมถึงจัดการฐานความรู้ที่ใช้สำหรับระบบ Chatbot ผ่านแดชบอร์ดผู้ดูแลระบบ
\end{enumerate}

ในด้านสถาปัตยกรรม ระบบได้รับการออกแบบเป็นบริการอิสระ 3 ส่วน (เว็บแอปพลิเคชัน บริการข้อมูลหลัก และบริการปัญญาประดิษฐ์) ที่สื่อสารระหว่างกันผ่าน RESTful API เพื่อให้สามารถพัฒนาและทดสอบแต่ละส่วนแยกจากกันได้ โดยองค์ประกอบเสริมสำหรับรองรับการขยายระบบในระดับใช้งานจริง เช่น API Gateway และ Message Broker จะได้รับการพัฒนาเพิ่มเติมในระยะถัดไปของโครงการ

ขอบเขตของข้อมูลในการวิจัยครอบคลุมจังหวัดขอนแก่น โดยในระยะต้นแบบเน้นการรวบรวมข้อมูลสถานที่ท่องเที่ยว ร้านอาหาร ที่พัก และกิจกรรมในเขตอำเภอเมืองขอนแก่นเป็นหลัก ร่วมกับอำเภอที่มีศักยภาพด้านการท่องเที่ยวที่คัดเลือกไว้เบื้องต้น (เช่น ภูเวียง อุบลรัตน์ ภูผาม่าน ชุมแพ น้ำพอง และหนองเรือ) และมีแผนขยายการเก็บข้อมูลให้ครอบคลุมอำเภออื่นภายในจังหวัดเพิ่มเติมในระยะถัดไปของงานวิจัย

การประเมินผลของระบบมุ่งเน้นด้านความถูกต้องและความเหมาะสมของแผนการเดินทาง ความถูกต้องของการตอบคำถามของระบบ Chatbot ประสิทธิภาพของการค้นคืนข้อมูลผ่าน RAG ระยะเวลาในการสร้างแผนการเดินทาง และความพึงพอใจของผู้ใช้งาน โดยประเมินผ่านการทดสอบเชิงประสิทธิภาพและแบบสอบถามจากผู้ใช้งานจริง

\subsection{ข้อจำกัดของโครงงาน}

\begin{enumerate}[label=(\arabic*)]
\item ข้อจำกัดด้านพื้นที่และข้อมูล การวิจัยจำกัดขอบเขตข้อมูลเฉพาะจังหวัดขอนแก่น โดยในระยะต้นแบบข้อมูลยังกระจุกตัวอยู่ในเขตอำเภอเมืองและอำเภอท่องเที่ยวสำคัญบางส่วนเท่านั้น ยังไม่ครอบคลุมทุกอำเภอของจังหวัด ส่งผลให้การนำระบบไปประยุกต์ใช้กับพื้นที่อื่นหรือขยายความครอบคลุมภายในจังหวัดจำเป็นต้องรวบรวมและจัดเตรียมฐานข้อมูลเพิ่มเติม นอกจากนี้ ความครบถ้วนและความถูกต้องของข้อมูลยังขึ้นอยู่กับแหล่งข้อมูลภายนอก (Google Places API) และข้อมูลที่ผู้ดูแลระบบนำเข้าสู่ระบบ
\item ข้อจำกัดด้านการอัปเดตข้อมูล ข้อมูลสถานที่ท่องเที่ยว ร้านอาหาร และกิจกรรมอาจมีการเปลี่ยนแปลงอยู่ตลอดเวลา เช่น เวลาเปิด--ปิด การย้ายสถานที่ หรือการยกเลิกกิจกรรม การนำเข้าข้อมูลในระยะต้นแบบดำเนินการเป็นครั้งคราวผ่านกระบวนการดึงและนำเข้าข้อมูล มิใช่การซิงค์ข้อมูลแบบเรียลไทม์ หากไม่มีการปรับปรุงข้อมูลอย่างสม่ำเสมอ อาจส่งผลให้คำแนะนำของระบบไม่สอดคล้องกับสถานการณ์จริง
\item ข้อจำกัดด้านเทคนิคของโมเดลปัญญาประดิษฐ์ แม้ระบบจะใช้สถาปัตยกรรม RAG ร่วมกับกลไกตรวจจับคำตอบที่ไม่มีข้อมูลรองรับ (Grounded Generation และ Hallucination Guard) เพื่อลดปัญหาการสร้างข้อมูลที่คลาดเคลื่อน (Hallucination) ของ LLMs แต่ยังคงมีโอกาสที่โมเดลจะสร้างคำตอบที่ไม่ถูกต้องหรือไม่สอดคล้องกับบริบทของข้อมูลได้ โดยเฉพาะในกรณีที่ฐานความรู้ไม่มีข้อมูลรองรับหรือข้อมูลไม่เพียงพอ
\item ข้อจำกัดด้านค่าใช้จ่ายในการดำเนินงาน ระบบมีต้นทุนในการให้บริการ เช่น ค่าใช้จ่ายของเซิร์ฟเวอร์ ค่าใช้บริการ API ของแบบจำลองภาษา และค่าใช้จ่ายในการดูแลฐานข้อมูลที่จัดเก็บทั้งข้อมูลทั่วไปและข้อมูลเวกเตอร์ ซึ่งอาจส่งผลต่อความสามารถในการขยายระบบหรือรองรับผู้ใช้งานจำนวนมากในอนาคต
\item ข้อจำกัดด้านสถาปัตยกรรมระบบ ระบบในระยะต้นแบบสื่อสารระหว่างบริการย่อยผ่าน RESTful API แบบ Synchronous เท่านั้น ยังไม่มีองค์ประกอบเสริม เช่น API Gateway หรือ Message Broker สำหรับการสื่อสารแบบ Asynchronous ซึ่งจำเป็นต่อการรองรับปริมาณผู้ใช้งานจำนวนมากในระดับใช้งานจริง องค์ประกอบดังกล่าวเป็นแนวทางการพัฒนาต่อยอดในระยะถัดไปของโครงการ
\item ข้อจำกัดด้านระบบยืนยันตัวตนผู้ใช้ ระบบสมัครสมาชิกและเข้าสู่ระบบในระยะต้นแบบยังอยู่ในขั้นตอนการพัฒนา โดยกลไกการยืนยันตัวตนและการจัดการสิทธิ์ผู้ใช้ (Authentication \& Authorization) ที่สมบูรณ์และมีความปลอดภัยในระดับใช้งานจริงจะได้รับการพัฒนาต่อในระยะถัดไปของโครงการ
\item ข้อจำกัดของระบบสะสมคะแนน (Point System) ระบบสะสมคะแนนจากการสแกน QR Code เป็นความสามารถเพิ่มเติมของระบบ ซึ่งประสิทธิภาพในการสร้างการมีส่วนร่วมของผู้ใช้งานขึ้นอยู่กับจำนวนสถานที่ที่เข้าร่วมโครงการ ความร่วมมือจากหน่วยงานที่เกี่ยวข้อง และการดำเนินกิจกรรมส่งเสริมการท่องเที่ยวของจังหวัดขอนแก่น ซึ่งอยู่นอกเหนือการควบคุมโดยตรงของทีมวิจัย
\end{enumerate}

% -----------------------------------------------------------
\section{ประโยชน์ที่คาดว่าจะได้รับ}
% -----------------------------------------------------------

\begin{enumerate}[label=\arabic*)]
    \item นักท่องเที่ยวสามารถเข้าถึงข้อมูลแหล่งท่องเที่ยว ร้านอาหาร และกิจกรรมในจังหวัดขอนแก่นได้อย่างครบถ้วนและสะดวกในที่เดียว
    \item ช่วยให้ผู้ใช้สามารถวางแผนทริปที่เหมาะสมกับความสนใจ เวลา และงบประมาณได้อย่างมีประสิทธิภาพมากขึ้น โดยระบบจะแนะนำสถานที่และกิจกรรมที่สอดคล้องกับความสนใจของผู้ใช้งาน
    \item ส่งเสริมการใช้ AI และ Chatbot ในการให้ข้อมูลเชิงโต้ตอบ ทำให้ประสบการณ์การท่องเที่ยวของผู้ใช้มีความเป็นส่วนตัวและทันสมัย
    \item หน่วยงานที่เกี่ยวข้องสามารถนำข้อมูลจากระบบไปใช้ประโยชน์ในการวิเคราะห์เชิงนโยบาย เช่น การจัดสรรทรัพยากรด้านการท่องเที่ยวหรือการวางแผนกลยุทธ์การตลาด
    \item เป็นกรณีศึกษาและต้นแบบในการพัฒนา Smart Tourism Platform ที่สามารถขยายผลไปยังจังหวัดอื่น ๆ หรือพัฒนาต่อยอดในระดับภูมิภาคได้
\end{enumerate}
